\documentclass[../main]{subfiles}
\begin{document}
\chapter{Carga eléctrica en un capacitor}
\img{images/08/capacitor}{Capacitores.}

\info{Carga almacenada en un capacitor}{
 La capacidad eléctrica, es la propiedad que tienen los cuerpos para mantener una carga eléctrica. La capacidad es también una medida de la cantidad de energía eléctrica almacenada para una diferencia de potencial eléctrico dada. El dispositivo más común que almacena energía de esta forma es el condensador. La relación entre la diferencia de potencial (o tensión) existente entre las placas del condensador y la carga eléctrica almacenada en este, se describe mediante la siguiente expresión matemática:
 \begin{equation*}
     Q=C\times V
 \end{equation*}
 Siendo
 
 $Q:$ La carga almacenada en Coulombs [C].
 
 $C:$ La capacidad en faradios [F].
 
 $V:$ El voltaje entre las placas del condensador en volts [V].
 
}



\actividad{Resolver los siguientes problemas:}
\begin{enumerate}
    \item Un condensador que se conecta a una tensión de $V=200V$, adquiere cada armadura una carga de $Q=6x10^{-9}C$. Calcular su capacidad.
    \item La carga que adquiere un condensador al conectarlo a $V=900V$ es de $Q=0,007C$. Calcular su capacidad.
    \item ¿A qué tensión debemos conectar un condensador de $C=104F$ para que adquiera una carga de $Q=2x10^{-6}C$?
    \item calcular la carga de un condensador de $C=200pF$ si se conecta a $V=100V$
    \item Un condensador de $C=2\mu F$ se conecta a $V=220V$. Calcular la carga. 
   \item Un condensador tiene una capacidad de 50 $\mu$F y se le aplica una tensión de 12 V. ¿Cuál es la \textbf{carga} almacenada en el condensador?
    \item Un condensador almacena una carga de 0.25 C cuando se le aplica una tensión de 25 V. ¿Cuál es la \textbf{capacidad} del condensador?
    \item Un condensador con una capacidad de 100 nF se carga hasta 10 V. ¿Cuál es la \textbf{carga} en el condensador?
    \item Un condensador se carga hasta 200 mC cuando se le aplica una tensión de 400 V. ¿Cuál es la \textbf{capacidad} del condensador?
    \item Si un condensador con una capacidad de 2 F almacena una carga de 10 C, ¿qué \textbf{tensión} se le ha aplicado?
    \item Un condensador tiene una capacidad de 470 $\mu$F y se le aplica una tensión de 5 V. ¿Cuál es la \textbf{carga} almacenada en el condensador?
    \item Un condensador de 1 mF se carga a 9 V. ¿Cuál es la \textbf{carga} en el condensador?
    \item Un condensador almacena una carga de 0.5 C cuando se le aplica una tensión de 50 V. ¿Cuál es la \textbf{capacidad} del condensador?
    \item Si un condensador con una capacidad de 10 nF se carga hasta 200 V, ¿cuál es la \textbf{carga} en el condensador?
    \item Un condensador tiene una capacidad de 330 $\mu$F y se le aplica una tensión de 6 V. ¿Cuál es la \textbf{carga} almacenada en el condensador?
    \item Un condensador almacena una carga de 0.1 C cuando se le aplica una tensión de 20 V. ¿Cuál es la \textbf{capacidad} del condensador?
    \item Un condensador con una capacidad de 47 nF se carga hasta 15 V. ¿Cuál es la \textbf{carga} en el condensador?
    \item Un condensador se carga hasta 2 mC cuando se le aplica una tensión de 500 V. ¿Cuál es la \textbf{capacidad} del condensador?
    \item Si un condensador con una capacidad de 3 F almacena una carga de 6 C, ¿qué \textbf{tensión} se le ha aplicado?
    \item Un condensador tiene una capacidad de 220 $\mu$F y se le aplica una tensión de 9 V. ¿Cuál es la \textbf{carga} almacenada en el condensador?
    \item Un condensador de 2 mF se carga a 12 V. ¿Cuál es la \textbf{carga} en el condensador?
    \item Un condensador almacena una carga de 1 C cuando se le aplica una tensión de 100 V. ¿Cuál es la \textbf{capacidad} del condensador?
    \item Si un condensador con una capacidad de 22 nF se carga hasta 300 V, ¿cuál es la \textbf{carga} en el condensador?
    \item Un condensador tiene una capacidad de 150 $\mu$F y se le aplica una tensión de 7 V. ¿Cuál es la \textbf{carga} almacenada en el condensador?
    \item Un condensador almacena una carga de 0.75 C cuando se le aplica una tensión de 75 V. ¿Cuál es la \textbf{capacidad} del condensador?
\end{enumerate}
\end{document}